\chapter{Linear Mappings}

\section{Definitions}
\begin{definitionenv}
	Let $V$ and $W$ be vector spaces over the same field $\mathcal{K}$.
	A mapping $T : V \to W$ is called a \textbf{linear mapping} (or \textbf{linear transformation}) if for all
	$\bm{u}, \bm{v} \in V$ and all scalars $\alpha, \beta \in \mathcal{K}$, we have
	\[
	T(\alpha \bm{u} + \beta \bm{v})
	= \alpha T(\bm{u}) + \beta T(\bm{v}).
	\]
\end{definitionenv}

\begin{exampleenv}
	Let $A \in \M_{m \times n}(\mathcal{K})$.
	Define a mapping
	\[
	T_A : \mathcal{K}^n \to \mathcal{K}^m, \quad
	T_A(\bm{x}) = A\bm{x},
	\]
	for all $\bm{x} \in \mathcal{K}^n$.
	Then $T_A$ is a \textbf{linear mapping}.

	Indeed, for any $\bm{x}, \bm{y} \in \mathcal{K}^n$ and $\alpha, \beta \in \mathcal{K}$,
	\[
	T_A(\alpha \bm{x} + \beta \bm{y})
	= A(\alpha \bm{x} + \beta \bm{y})
	= \alpha A\bm{x} + \beta A\bm{y}
	= \alpha T_A(\bm{x}) + \beta T_A(\bm{y}).
	\]
	Hence $T_A$ satisfies the linearity condition,
	and matrix left-multiplication defines a linear mapping from $\mathcal{K}^n$ to $\mathcal{K}^m$.
\end{exampleenv}

\begin{exampleenv}
	Let $V = \mathcal{K}^n$ and $W = \mathcal{K}^m$.
	Denote by $\mathcal{L}(V, W)$ the set of all linear mappings from $V$ to $W$ over the field $\mathcal{K}$.
	That is,
	\[
	\mathcal{L}(V, W)
	= \{\, T : V \to W \mid T \text{ is linear} \,\}.
	\]

	For $T, F \in \mathcal{L}(V, W)$ and $\alpha \in \mathcal{K}$,
	we define their \textbf{sum} and \textbf{scalar multiple} by
	\[
	(T + F)(\bm{u}) = T(\bm{u}) + F(\bm{u}),
	\qquad
	(\alpha T)(\bm{u}) = \alpha\, T(\bm{u}),
	\quad \forall\, \bm{u} \in V.
	\]

	With these operations of addition and scalar multiplication,
	$\mathcal{L}(V, W)$ forms a vector space over $\mathcal{K}$.
\end{exampleenv}


\begin{definitionenv}
	Let $V$ be a vector space over a field $\mathcal{K}$, and let $\bm{a} \in V$.
	The mapping $T_{\bm{a}} : V \to V$ defined by
	\[
	T_{\bm{a}}(\bm{x}) = \bm{x} + \bm{a},
	\]
	for all $\bm{x} \in V$, is called a \textbf{translation} by the vector $\bm{a}$.

	A translation shifts every vector in $V$ by the same displacement $\bm{a}$.
	Note that a translation is generally \textbf{not} a linear mapping, since
	\[
	T_{\bm{a}}(\bm{0}) = \bm{a} \neq \bm{0}.
	\]
\end{definitionenv}

\begin{theoremenv}
	Let $V$ and $W$ be vector spaces over a field $\mathcal{K}$.
	Let $\{\bm{v}_1, \dots, \bm{v}_n\}$ be a basis of $V$,
	and let $\bm{w}_1, \dots, \bm{w}_n$ be arbitrary elements of $W$.
	Then there exists a unique linear mapping
	\[
	T : V \to W
	\]
	such that
	\[
	T(\bm{v}_i) = \bm{w}_i, \quad i = 1, \dots, n.
	\]

	Moreover, if $x_1, \dots, x_n \in \mathcal{K}$, then
	\[
	T(x_1 \bm{v}_1 + x_2 \bm{v}_2 + \cdots + x_n \bm{v}_n) = x_1 \bm{w}_1 + x_2 \bm{w}_2 + \cdots + x_n \bm{w}_n.
	\]
\end{theoremenv}

\section{Kernel and Image}
\begin{definitionenv}
	Let $T : V \to W$ be a linear mapping between vector spaces over a field $\mathcal{K}$.
	The \textbf{kernel} of $T$, denoted by $\ker(T)$, is the subset of $V$ defined by
	\[
	\ker(T) = \{\, \bm{v} \in V \mid T(\bm{v}) = \bm{0} \,\}.
	\]
	It consists of all vectors in $V$ that are mapped to the zero vector in $W$.

\end{definitionenv}

\begin{remarkenv}
	The kernel $\ker(T)$ is a subspace of $V$.
\end{remarkenv}

\begin{theoremenv}
	Let $F : V \to W$ be a linear map between vector spaces over a field $\mathcal{K}$.
	The following statements are equivalent:
	\begin{enumerate}
		\item $\ker(F) = \{\bm{0}\}$;
		\item For all $\bm{v}, \bm{w} \in V$, if $F(\bm{v}) = F(\bm{w})$, then $\bm{v} = \bm{w}$.
	\end{enumerate}
	In other words, $F$ is injective if and only if its kernel is $\{\bm{0}\}$.
\end{theoremenv}

\begin{remarkenv}
	Note that $F$ is a mapping if and only if $F^{-1}$ is injective and surjective.
\end{remarkenv}

\begin{theoremenv}
	Let $F : V \to W$ be a linear map whose kernel is $\{\bm{0}\}$.
	If $\bm{v}_1, \dots, \bm{v}_n$ are linearly independent elements of $V$,
	then $F(\bm{v}_1), \dots, F(\bm{v}_n)$ are linearly independent in $W$.
\end{theoremenv}

\begin{definitionenv}
	Let $T : V \to W$ be a linear map between vector spaces over a field $\mathcal{K}$.
	The \textbf{image} of $T$, denoted by $\operatorname{im}(T)$, is the subset of $W$ defined by
	\[
	\operatorname{im}(T) = \{\, T(\bm{v}) \mid \bm{v} \in V \,\}.
	\]
	It consists of all vectors in $W$ that can be expressed as the image of some vector in $V$ under $T$.

\end{definitionenv}

\begin{remarkenv}
	The image $\operatorname{im}(T)$ is a subspace of $W$.
\end{remarkenv}

\begin{theoremenv}[Rank--Nullity Theorem]
	Let $T : V \to W$ be a linear map between finite-dimensional vector spaces over $\mathcal{K}$.
	Then
	\[
	\dim(\ker T) + \dim(\operatorname{im} T) = \dim V.
	\]
	Here, $\dim(\ker T)$ is called the \textbf{nullity} of $T$,
	and $\dim(\operatorname{im} T)$ is called the \textbf{rank} of $T$.
\end{theoremenv}

\begin{theoremenv}[Rank--Nullity Theorem for Matrices]
	Let $A \in \M_{m \times n}(\mathcal{K})$, and let its column vectors be
	\[
	\bm{A}_1, \bm{A}_2, \dots, \bm{A}_n \in \mathcal{K}^m.
	\]
	We call
	\[
	\operatorname{Col}(A) = \operatorname{span}\{\bm{A}_1, \dots, \bm{A}_n\}
	\]
	the \textbf{column space} of $A$, and
	\[
	\operatorname{Null}(A) = \{\bm{x} \in \mathcal{K}^n \mid A\bm{x} = \bm{0}\}
	\]
	the \textbf{null space} of $A$.

	Then $\operatorname{Col}(A)$ is a subspace of $\mathcal{K}^m$,
	and $\operatorname{Null}(A)$ is a subspace of $\mathcal{K}^n$.
	Moreover,
	\[
	\dim(\operatorname{Col}(A)) + \dim(\operatorname{Null}(A)) = n,
	\]
	where $n$ is the number of columns of $A$.
\end{theoremenv}


\begin{theoremenv}
	Let $L : V \to W$ be a linear map between finite-dimensional vector spaces over $\mathcal{K}$.
	Assume that
	\[
	\dim V = \dim W.
	\]
	If $\ker L = \{\bm{0}\}$, or if $\operatorname{im} L = W$,
	then $L$ is \textbf{bijective}.
\end{theoremenv}

\begin{definitionenv}
	If $f$ is a mapping, and its inverse correspondence is also a mapping, then we say $f$ is \textbf{invertible}.
\end{definitionenv}

\begin{propositionenv}
	A mapping $f$ is invertible if and only if it is bijective.
\end{propositionenv}

\begin{theoremenv}
	Suppose $V$ and $W$ are vector spaces over a field $\mathcal{K}$ and $T: V \to W$ is a linear map that is invertible. Then its inverse $T^{-1}: W \to V$ is also a linear map.
\end{theoremenv}

\begin{theoremenv}
	Let $V$ and $W$ be vector spaces over a field $\mathcal{K}$ with
	\[\dim V = \dim W,\]
	Suppose a linear map $T:V \to W$, if $\ker T = \{\bm{0}\}$, then $T$ is invertible.
\end{theoremenv}

\section{Isomorphism}
\begin{definitionenv}
	Suppose $V$ and $W$ are vector spaces over a field $\mathcal{K}$. Any invertible linear map $T: V \to W$ is called an \textbf{isomorphism} between $V$ and $W$.
\end{definitionenv}
\begin{notationenv}
	We say that $V$ and $W$ are \textbf{isomorphic} if there exists an isomorphism $T : V \to W$. In this case we denote this fact by
	\[
	V \cong W .
	\]
\end{notationenv}
\begin{theoremenv}
	Let $V$ and $W$ be finite-dimensional vector spaces over a field $\mathcal{K}$. Then $V \cong W$ if and only if $\dim V = \dim W$.
\end{theoremenv}

\begin{remarkenv}
	Suppose $V$ and $W$ are vector spaces over a field $\mathcal{K}$ with
	\[
	\dim V = m, \quad \dim W = n.
	\]
	Then $V \cong \mathcal{K}^m$ and $W \cong \mathcal{K}^n$.
\end{remarkenv}

\begin{theoremenv}
	Let $U,\, V,\, W$ be vector spaces over a field $\mathcal{K}$, with linear maps $F: U \to V$ and $G: V \to W$. Then the composition map \[G \circ F: U \to W\] is also a linear map.
\end{theoremenv}


\newpage